\section{Context}
Outlier Analysis is an important task for research and real world applications.
In recent years, machine learning is increasingly used for outlier detection tasks.
\citep{shenkar2022anomaly} recently published a interesting machine learning approach based on contrastive learning, which I will focus on.
Oulier detection at its core is an classification task with highly imbalanced data.
Nevertheless, a uniform definition of outliers could yet not be found. In 2017 \citet{ayadi2017outlier} found 12 different interpretations of outlier.
All of them dont differ so much, so I will use the definition of Hawkins, because of the influece his paper had on this topic.
Ouliers, also called anomalies \citep{chandola2009anomaly}, are observations, that deviate extremely from the defined norm.
They are the product of a different mechanisms, that aren't included into the research question. \citep{hawkins1980identification}
On the one hand, depending on the context they can be researchers target and provide useful information.
For instance outliers develope by measuring recording errors, misreporting or sample errors.
On the other hand, if they are not the target, they have to be identified and carefully treated. \citep{smiti2020critical}
For more information about outlier treatment I recommend.
\todo{insert survey about outlier treatment}

High stakes real world application can be found, where outlier detection plays a crucial role.
In Cybersecurity, outlier detection can be used to detect intrusions or maleware throuigh unusual network traffic.
\todo{insert paper that this describes}
In financial domains, outliers are been used to identify criminal activity.
For example in fraudulent financial behaviour.
The difference in usage of a stolen credit card by the criminal, can be detected by outlier detection methods. \citep{rezapour2019anomaly}
In Healthcare, outlier detection methods can be used to detect early rare diseases or help medical staff by diagnosing. 


There are 3 different types of outliers: 
\begin{enumerate}
\item Global outliers, sometimes named point anomalies, are the simplest type of anomaly. 
When a obervation with respect to rest of the its distribution has extreme values, you could consider it a outlier.
E.g. a Person spends 100€ a week and someday the person spends 1000€. that would be possibly be a global outlier, because it is 
a extreme value out of the normal distribution. 
\item Contextual outliers, sometimes named conditional outliers, are outliers that are outliers bescause of their specific contexts.
A data instance is defined by it behavioral and contextual attributes. 
Contextual attributes define contextual atrributes and behavioral attributes define non-contextual attributes.
Using the credit card fraud example, behavioral atrributes are the amount of paymets done in a certain period.
Contextual attribute is the time dimension of this longitudinal data. 
If a person spends 100€ in a week, accept for christmas holiday. There the Person spends 1000€.
the contextual outlier would be, if a transaction of 1000€ emerges in mid july.
1000€ is not an extreme charakteristic, because it will be spend in a different time.
But it is a unusual period of time. 
\item Collective outliers are observations that are by itselfe not unsual but their accurance together as a collection. 
Using the example as before, a single use of a credit card for a normal payment is not a outlier, but when a person pays 50 times
a small payment, that would be considered a collective outlier.
\end{enumerate}


There are different other types of detection methods, that have to be circumvented from outlier detection.
Noisy data has no useful information and hinders the analysis. 
It need to be identified and corrected because it makes computation difficult or impossible to do.
Noisy data contains incorrect data types, missing values or incorrect data \citep{smiti2020critical}.


Novelty Detection, often synonymously used with outlier detection, targets new and previously unconsidered patterns in the data \citep{domingues2019probabilistic}.
The findings are normally incorporated into a new model and not like in outlier detection being removed or eximaned further \citep{chandola2009anomaly}.

Out-Of-Destribution detection (OOD) targets unfamiliar datapoints.
Algorithms classify datapoint in different classes and the goal is find datapoints that wont fit into the classes.
Openset recognition is closely related to OOD. 
The classification task differs from OOD, such as the algorithm has classify known categories correctly while rejecting unknown classes.


A fairly new approach in outlier detection is to provide meaningful explanations for outliers.
There is a need to tell not just wich observation is a outlier, but to tell why that specific outlier is a outlier.
E.g. in Fraud detection, analyst need to understand why a transaction was being flagged as fraudulent and wether it need human intervention.
A common approach is to a ssign a importnace level to the outlier.
The more far away an outlier is from the "normal" data distribution, the more likely it need human attention 
Therefor a user can prioritize. 


A different type a providing meaningfull insights in outliers to highlight causal interactions among outliers. 
Causal in a sense, that one outlier produced other outliers in chronological way.
So if you remove that first outlier, you would be removing the other outliers too.
For example a traffic jam in one part of the city is beeing caused of another part of the city
Users therefor can use that information to prevent the first outlier, so that other outliers don't develope \citep{panjei2022survey}.



\subsubsection*{Challenges}
The curse of dimenionality (CoD) plays a big role in highdimensional datasets.
However, in the big data era, the sheer number of variables that can be collected from a single sample can be
problematic. we discuss four important problems of dimensionality as it applies to data sparsity, multicollinearity, testing and overfitting
First, as the dimensionality p increases, the volume that the samples may occupy grows rapidly. If we treat the distance between points
(e.g., Euclidian distance) as a measure of similarity, then we interpret greater distance as greater dissimilarity. As p increases, this
dissimilarity increases because the mean distance between points increases as wurzelp (Fig. 2a). This effect is stark at high values of
p. For example, at p = 100, the closest pair of points are farther from one another than the most distant two points are for about p < 15
400this month The assessment of whether a particular data value is an outlier accordingly also becomes difficult. For example, 68\% of normally
distributed points in one dimension fall within sigma of the mean, but this fraction drops off precipitously for large values of p (Fig. 2b).
For example, at p = 10, the fraction of points within sigma of the mean is only 0.017\%, which is equivalent to the points within 0.00022sigma in
one dimension (Fig. 2b, dashed lines). Putting it another way, points within sigmaat p = 10 are as rare as points outside of 3.8sigma at p = 1. 
These are not intuitive observations—proportions of how points are distributed at higher dimensions are not the same as in one dimension.
hat if we use a different distance measure to express similarity between subjects, such as correlation? Unfortunately, we cannot escape the CoD—the range of
correlations between points drops rapidly (Fig. 2c) with p. For example, at p = 100, among 10,000 random pairs of points we see no correlations greater than 0.5, and most
are extremely tightly grouped around 0. As the number of variables increases, the number of subjects in each set of categories decreases and the
correlation between any two subjects across variables also decreases.
When the number of dimensions is larger than the number of samples (p > n), another effect that confounds analysis is perfect multicollinearity.
In this case, we can always express at least one of the variables as a linear combination of the others. For example, if we
have p = 3 variables but only n = 2 subjects and we think of the subjects as two 3D vectors, then from linear algebra we know that
these vectors define either a point (i.e., they have the same three values) or a line. In both cases, the three variables are related linearly.
This is the case any time there are fewer samples than dimensions—the variables span a lower-dimensional subspace in which some
of the dimensions become redundant and expressible in terms of other dimensions, thus yielding perfect multiple correlation.
Finally, overfitting is another CoD that occurs because the flexibility of prediction equations5 is in part determined
by the number of variables involved. With increased flexibility, prediction and classification rules adapt to both the
patterns in the population and the random idiosyncrasies of the training sample. \citep[p.1pp]{altman2018curse}
\todo{research: embeddings to solve the curse of dimensionality}


scarcity of Lables

Concept drift


A paper published by \citep{shenkar2022anomaly} introduced a newly method, based on a internal contrastive learning approach, called InterCont.
The authors compared their methods agaist a wast number of benchmark methods and supllied comprehensive results.
Their method achieved high values on their chosen metrics. 
Surprisingly the well established k-NN algorithm achieved a higher mean and a lower variance, especially in the F1 metric compared to the InterCont.


I will take this result to establish my own set of experiments and compare 3 different outlier detection methods. 
The InterCont will be the primary method I would like to analyse. 
As an classic machine learning benchmark method and to analyse the result I mentioned above, I will include k-NN as well.
As a third method I will include the minimum covariance determinant (MCD) as a statistical benchmark.


I suppose 2 different Hypothesis in my work:
\begin{enumerate}
    \item H1: The k-NN will produce the best outcome of all methods. 
    \item H2: InterCont will produce the second best outcome. 
\end{enumerate}


My thesis will be structured as follows:
I begin with describing the theoretical basis for InterCont. 
Therefore I will concentrate on different machine learning paradigms and after that I will describe the basic arcitecture of machine learning algorithm.
In the next step I will describe contrastive learning and apply the arcitectual specifics of contrastive learning algorithms.
That will lead me to the description of the details in the InterCont algorithm.
After that I will transition to the benchmark methods, beginning with k-NN.
I will describe the statistical basis for MCD afterwards, compleating the theoretical basis for all methods.
After that I will transition to the implementation section.
There I will describe my chosen dataset and I will analyse with different visualisation techniques the caracteristics of my dataset.
After describing the experimental design and setup I will transiion to the results.
In the end I will apply the theory to the result and I will explain how the results fit into it. 
 